\domainevide{Mathlib contient les puissances divisées sur un anneau, mais pas le site cristallin ni la cohomologie de de Rham des schémas.}{4, 5, 6, 7, 8, 54, 137, 156-4}
\noindent Pistes de l'onglet Findings pour ces cotes, non démontrées ici :
\begin{itemize}\small\setlength{\itemsep}{0pt}
\item Cote~4 : si l'on admet le théorème de Lefschetz faible cristallin avec torsion, le nombre de valeurs propres unités du Frobenius sur la cohomologie cristalline évanescente d'une section hyperplane lisse est au moins le rang $F$-semi-simple de la cohomologie cohérente évanescente. Mathlib n'a pas la cohomologie cristalline.\piste{4-vanishing-unit-root-bound}
\item Cote~5 : la $W(k)$-algèbre locale artinienne universelle dont l'idéal maximal porte des puissances divisées nulles en degrés $\geq p^{\nu+1}$ est $W[\Lambda_0, \ldots, \Lambda_\nu]/(\Lambda_i^p - q_i\Lambda_{i+1}, \Lambda_\nu^p, \Lambda_0 - p)$, avec une forme normale explicite. Accessible en partie : l'intégralité des $q_i$ et leur valuation $p$-adique ; la forme normale demande des calculs dans les vecteurs de Witt.\piste{5-universal-artinian-pd-normal-form}
\item Cote~5 : la lettre du 10~décembre~1965 énonce la dualité étale pour un morphisme lisse compactifiable de sites annelés, pour des faisceaux d'anneaux quelconques tués par un entier premier aux caractéristiques résiduelles. Mathlib n'a pas la dualité étale.\piste{5-deligne-letter-ringed-duality}
\item Cote~6 : pour une polarisation de degré quelconque, divisible par $p$ compris, l'anneau local d'un point du schéma de modules des variétés abéliennes ordinaires polarisées est d'intersection complète. Mathlib n'a pas les schémas de modules de variétés abéliennes.\piste{6-ordinary-moduli-complete-intersection}
\item Cote~6 : sur un corps fini, la torsion du groupe formel d'une courbe elliptique supersingulière par une unité de norme réduite transcendante donne une forme de $G_{1,1}$ qui ne se plonge dans le groupe formel d'aucune variété abélienne. Mathlib n'a pas les groupes formels des variétés abéliennes ni les conjectures de Weil.\piste{6-formal-group-not-in-abelian-variety}
\item Cote~6 : sur un corps algébriquement clos de caractéristique $p$, il existe des surfaces abéliennes simples de même groupe formel $\hat{\mathbf G}_m \times G_{1,1}$ qui ne sont pas isogènes. Mathlib n'a pas les surfaces abéliennes ni la famille de Manin.\piste{6-same-formal-group-not-isogenous}
\item Cote~7 : le dossier propose comme cohomologie $p$-adique celle d'un site de schémas formels affines $p$-adiques, sans puissances divisées ni surconvergence, soumise à deux tests (Washnitzer--Monsky, de Rham d'un relèvement). Mathlib n'a pas de cohomologie $p$-adique des schémas.\piste{7-noninfinitesimal-site}
\item Cote~8 : quand $\pi A$ a des puissances divisées, le relèvement de Lubin--Tate existe et est unique par la seule théorie des déformations cristallines ; sans puissances divisées, l'argument donne l'unicité à isogénie près. Mathlib n'a pas la théorie des déformations cristallines.\piste{8-lubin-tate-via-hodge-filtration}
\item Cote~8 : deux groupes de Barsotti--Tate de dimension un et de hauteur $[K : \mathbf Q_p]$ sur $\mathbf F_q$, munis d'une action de $A$, sont isogènes si et seulement si $\varpi = \varpi'$, et $\mathrm{Hom}_A(G_0, G'_0)$ est alors explicite. Mathlib n'a pas les groupes de Barsotti--Tate.\piste{8-dim-one-A-module-homs}
\item Cote~137 : pour la catégorie $\mathcal P$ des complexes bornés de modules libres de type fini, l'évaluation sur le complexe des disques est une équivalence entre les foncteurs $\mathcal P \to \mathcal B$ exacts à gauche sur $\Sigma$ et les complexes de cochaînes de $\mathcal B$. Accessible avec les complexes homologiques de mathlib ; la preuve est longue.\piste{137-perfect-complexes-universal-sigma-exact}
\item Cote~137 : énoncé sans démonstration : les catégories polynomiales strictes sont équivalentes aux $\otimes$-pd-catégories. Mathlib n'a aucune des deux notions.\piste{137-polynomial-categories-pd-tensor-categories}
\end{itemize}
